\chapter{Background and Related Work}
\label{ch:background}

This thesis sits where four lines of work meet: keystroke biometrics, which
asks \emph{who} is typing; sequential change detection, which asks
\emph{when} something changed; the analysis of programming processes in
computing education, which records \emph{how} students write code; and
academic-integrity practice, which decides what may be done with any of it.
This chapter reviews each in turn, then the two concerns that limit every
behavioural measure of students, fairness and privacy, and closes with the gap
that the research questions of \cref{ch:problem} address.

\section{Keystroke dynamics}
\label{sec:rw-keystroke}

\subsection{Static and free-text authentication}
Keystroke dynamics identifies people by the timing of their typing: how long
each key is held (dwell or hold time) and how long passes between successive
keys (flight time, or digraph latency for a pair of keys). Surveys describe
decades of work and a wide range of features and classifiers
\cite{banerjee2012survey,teh2013survey}. The early systems were \emph{static}:
the user types a fixed string such as a password, and the timing vector is
compared with an enrolled template. Monrose and Rubin showed that typing
rhythm is distinctive enough to serve as a biometric for both structured and
free text \cite{monrose2000keystroke}. Bergadano et al.\ compared samples by
the relative order of trigraph durations rather than their absolute values,
which made the measure robust to a typist who is faster or slower on a given
day \cite{bergadano2002keystroke}.

Free-text methods drop the fixed string. Gunetti and Picardi compared any two
samples through the n-graphs they share, using a relative (rank-order) and an
absolute distance, and showed that people can be told apart from ordinary
typing of different texts \cite{gunetti2005freetext}. This matters here,
because students in a lab never type the same text twice: every question, and
every student's solution to it, is different.

Benchmarks made the field comparable. Killourhy and Maxion collected one
password typed 400 times by each of 51 people and compared 14 anomaly
detectors under one protocol; the scaled Manhattan detector, which divides
each feature's distance by its mean absolute deviation in the user's
training data, had the lowest equal error rate (9.6\%)
\cite{killourhy2009comparing}. Its simplicity and strength make it the
standard baseline, and it is one of the baselines of this thesis
(\cref{ch:rq1}). Deep models came later. TypeNet trains a recurrent network
with a triplet loss to embed sequences of key timings, so that a new user is
enrolled from a few samples without retraining, and evaluates on very large
populations \cite{acien2022typenet}. It was trained on the Aalto typing data of
Dhakal et al., 136 million keystrokes from 168,000 volunteers transcribing
sentences online, whose analysis also showed how much typing varies with
skill and habit (for example, rollover, pressing the next key before releasing
the current one) \cite{dhakal2018typing}. A TypeNet-style embedding is the
second learned baseline in this thesis.

\subsection{Continuous authentication}
Authentication at log-in says nothing about who types afterwards.
\emph{Continuous} authentication keeps checking during a session, typically
over a sliding window of recent keystrokes. Bours argued that such systems
should be judged not by an equal error rate on fixed-size samples but by how
many actions an impostor can take before being locked out and how many a
genuine user takes before a false lock-out, and proposed a trust model in
which each action raises or lowers a running trust score
\cite{bours2012continuous}. Mondal and Bours combined keystroke and mouse
dynamics in this framework \cite{mondal2017continuous}. On phones, continuous
authentication from touch, motion and typing is an active area of its own
\cite{patel2016continuous}.

Shared free-text data sets let these ideas be tested on natural typing. The
Buffalo data set collected long free-text and transcription sessions,
including sessions on a different keyboard \cite{sun2016buffalo}; the
Clarkson~II data set logged participants' ordinary computer use, unprompted,
over a long period \cite{murphy2017clarkson}. Both report that accuracy falls
when typing is natural rather than prompted and when hardware changes, which
is the setting of this thesis.

\subsection{Evaluating behavioural biometrics}
Two critiques shape the evaluation in \cref{ch:study}. Eberz et al.\ showed
that common shortcuts inflate reported accuracy: training and testing on
interleaved rather than time-ordered samples, ignoring how a user's behaviour
drifts between sessions, and summarising a system by one equal error rate
\cite{eberz2017evaluating}. Sugrim et al.\ likewise argued that a single
error rate hides how errors are spread across users and thresholds, and
proposed metrics that expose it \cite{sugrim2019metrics}. This thesis
therefore splits data by student and by week, always scores later sessions
against earlier baselines, reports false alarms per hour and per session
rather than an equal error rate, and gives subgroup error rates (\cref{ch:rq4}).

\subsection{From authentication to authorship continuity}
Continuous authentication asks ``is this the enrolled user?'' at every
window. The question of this thesis is narrower and, in one way, easier: has
the person typing \emph{changed} since the session began? The answer does not
need an absolute decision per window, only evidence that accumulates after a
change. It is also harder in another way: the typing is source code, full of
navigation, deletion and long pauses for thought, produced by novices whose
typing is still developing, on whatever device the lab offers. These
differences motivate the change-point formulation below.

\section{Change-point detection}
\label{sec:rw-changepoint}
Sequential change-point detection watches a stream of observations and raises
an alarm as soon as possible after their distribution changes, while keeping
false alarms rare. Page's cumulative sum (CUSUM) procedure accumulates the
log-likelihood ratio of ``after'' against ``before'', resetting at zero, and
alarms when the sum crosses a threshold $h$ \cite{page1954continuous}. Lorden
framed the problem as minimising the worst-case expected delay subject to a
bound on the mean time between false alarms, and showed CUSUM to be
asymptotically optimal under that criterion \cite{lorden1971procedures};
Moustakides proved it exactly optimal \cite{moustakides1986optimal}. The
standard references cover the theory and its many variants
\cite{basseville1993detection,tartakovsky2014sequential}. In statistical process
control the same procedure is the CUSUM chart, tuned by a reference value $k$
(the smallest shift worth detecting, in standard deviations) and a decision
interval $h$ \cite{hawkins1998cusum}; the detector of \cref{ch:detector} uses
exactly these two parameters per channel, and the classroom results
(\cref{sec:retune}) are largely a story about choosing $k$.

Two properties of CUSUM suit authorship continuity. It needs only the
``before'' distribution, which a per-student baseline supplies, and a
size of shift, not a model of every possible impostor. And the time at which
the sum last left zero is a natural estimate of when the change happened,
which the review timeline shows to the teacher. Offline methods, which look
at a whole recorded series and segment it after the fact, are surveyed by
Aminikhanghahi and Cook \cite{aminikhanghahi2017survey} and by Truong et al.\
\cite{truong2020selective}; Bayesian online change-point detection maintains a
posterior over the time since the last change \cite{adams2007bocpd}. These
could re-segment a finished session for review, but a live classroom needs a
decision while the student is still there, and CUSUM's two parameters can be
explained to a teacher.

\section{Programming-process data in computing education}
\label{sec:rw-process}

\subsection{Learning analytics from the IDE}
Computing education has long logged how students program, not only what they
submit. Jadud's error quotient summarised sequences of compilations by how
often the same error recurred \cite{jadud2006methods}. Blikstein et al.\
clustered snapshots of students' code over time and found distinct, stable
programming styles \cite{blikstein2014pluralism}. An ITiCSE working group
reviewed educational data mining in programming and found that most work
predicts performance or detects struggle, from data at the granularity of
compilations or submissions \cite{ihantola2015edm}. Hundhausen et al.\ proposed
a process model for IDE-based learning analytics, from collecting data to
intervening, and called for more attention to how analytics are acted on
\cite{hundhausen2017ide}. ProgSnap2 standardised the format of programming
process data so that data sets and tools can be shared
\cite{price2020progsnap2}. The event schema of this thesis
(\cref{app:schema}) is finer than ProgSnap2's typical granularity
(individual keys and edits with their origin) but follows the same idea of a
documented, versioned format.

\subsection{Keystrokes of programmers}
The closest strand uses keystroke timing in programming courses. Longi et al.\
identified students from the timing of their typing while programming, with
high accuracy once enough keystrokes were available
\cite{longi2015identification}. Leinonen et al.\ inferred programming
experience and course performance from the same kind of data
\cite{leinonen2016inference}, and studied whether typing patterns collected
during a course can authenticate students in practical programming exams
\cite{leinonen2016exams}. Edwards et al.\ compared keystroke data from writing
in a natural language and in a programming language, and found that the
context changes which features predict learning outcomes
\cite{edwards2020keystroke}. Hellas et al.\ followed up suspicious take-home
exam submissions and found a spectrum from help-seeking and collaboration to
systematic cheating \cite{hellas2017plagiarism}. Together these establish that
programmers' typing is individual and measurable in real courses, and that
integrity problems in programming are not only copying of finished code. None
of them watches for a change of typist \emph{within} a session.

\subsection{Keystroke logging in writing research}
Writing research has studied keystroke logs for longer. Tools such as Inputlog
record every key and pause of a writer and visualise the writing process
\cite{leijten2013inputlog}. Conijn et al.\ showed that the writing task itself
changes keystroke features (pause and revision behaviour differ between
copying and composing, and between kinds of composition)
\cite{conijn2019keystroke}. Crossley et al.\ used such process features to
separate original essays from transcribed ones with high accuracy, because
transcription is more linear and bursty \cite{crossley2024plagiarism}. The task
effect found by Conijn et al.\ reappears in the pilot of this thesis: build
order and revision follow the question as much as the person
(\cref{sec:retune}).

\subsection{Nosi IDE: the closest prior work}
Nosi \cite{lueder2026nosi} is a fork of VS Code with LLM features disabled.
Assignments are distributed as encrypted files that are decrypted only in
memory, so copying code out produces an encrypted payload and every paste
is logged with its timestamp and content. This leaves manual retyping
(``scribing'') of LLM output as the remaining stealthy route, which Nosi
targets with one feature: the L-scale (an L-moment measure of dispersion
\cite{hosking1990lmoments}) of log-transformed inter-arrival times (IAT)
between keystrokes, after removing pauses over 5 minutes and zero-length
intervals. Authentic coding has a long-tailed IAT distribution; scribing is
flatter.

The evidence is a proof-of-concept with $N = 5$, all authors and proficient
coders, in an A--B--B--A design (code, scribe, scribe, code) on one
priority-queue task. A linear SVM on the single feature reached 100\%
leave-one-participant-out accuracy on whole sessions and about 80\% on
\SI{150}{\second} slices; typing comments was sometimes misread as
scribing. In deployment in a course of about 1000 students, staff found 48
cases of LLM-assisted plagiarism by replaying edit logs.

Nosi differs from this thesis in three ways. It asks whether code was
\emph{transcribed}, not whether the \emph{author changed}; it uses a
population-level, subject-independent classifier rather than a per-student
baseline; and its sample is five expert authors. \RQ{2} relates its measure
to observable copying signs at larger $N$ with ordinary students, and \RQ{1}
addresses the switch it does not model.

\section{Academic integrity in programming assessment}
\label{sec:rw-integrity}

\subsection{Contract cheating and impersonation}
\emph{Contract cheating} is having a third party complete assessed work,
whether paid or not \cite{lancaster2016contract}. Newton's systematic review
of self-report surveys found that the share of students admitting to paying
someone for their work rose over time, to about one in seven in surveys from
2014 onward \cite{newton2018contract}. In a survey of over 14,000 students at
Australian universities, Bretag et al.\ found about 6\% admitting to at least
one contract-cheating behaviour, including having someone else sit an exam
\cite{bretag2019contract}. Self-reports undercount, and the rates differ by
country and discipline, so these numbers bound the problem only loosely; the
precision analysis of \cref{ch:rq4} therefore uses two prevalences, 1 in 100
and 1 in 500, rather than one estimate. Dawson treats the problem as one of
\emph{assessment security}: no measure stops a determined cheat, and what
matters is how much each measure raises the cost of cheating and how it fails
\cite{dawson2020defending}. A takeover partway through a lab is a cheap form
of contract cheating: no money changes hands, and the helper needs only the
student's seat or screen for a few minutes.

\subsection{Checks at the start: identity and proctoring}
Online and lab exams verify identity when a session starts, through a
log-in, an ID card or a face match, and remote proctoring adds webcam,
microphone and screen recording for the rest of the session. Atoum et al.\
built an automated proctor from a webcam, a wearable camera and a microphone
\cite{atoum2017proctoring}; Nigam et al.\ reviewed the AI-based proctoring
systems that spread during the COVID-19 pandemic \cite{nigam2021proctoring}.
These systems watch the room and the face, not the work, and they carry high
costs in privacy and trust: students describe them as invasive
\cite{balash2021examining}, ethicists question whether their benefits justify
the surveillance \cite{coghlan2021proctor}, and universities adopted them as
a ``necessary evil'' under emergency conditions \cite{selwyn2023necessary}.
A lab already has a teacher in the room; what the teacher cannot see is whose
hands produced which part of the code on sixty screens.

\subsection{Checks at the end: similarity and generated code}
Source-code plagiarism detection compares finished submissions. Winnowing
fingerprints \cite{schleimer2003winnowing}, greedy string tiling
\cite{wise1993gst} and tools such as JPlag \cite{prechelt2002jplag} find
shared structure that survives renaming and reordering; reviews cover the
many tools used in academia \cite{novak2019similarity} and the wider
literature on plagiarism in programming assessments \cite{albluwi2019plagiarism}.
Similarity detection catches two students who submit the same code. It cannot
see a helper who writes a student's only copy, and it says nothing about
\emph{when} or \emph{by whom} code was written.

\subsection{Checks of the process: pastes and edit logs}
Recording the process closes part of that gap. Editors used for assessment
increasingly log pastes and replay edit histories, as Nosi does; the editor of
this thesis logs every paste with its origin and blocks text pasted from
outside the editor when the teacher turns that on (\cref{sec:paste-guard}). A
paste log shows \emph{what} arrived without being typed. A takeover, however,
is typed: the helper's keystrokes look like any other keystrokes unless they
are compared with the student's own.

\section{Large language models and assessment}
\label{sec:rw-llm}
Code-generating language models changed what an unsupervised programming
task measures. Codex solved a large share of hand-written programming
problems from their docstrings \cite{chen2021codex}, and AlphaCode reached
roughly median human performance in programming competitions
\cite{li2022alphacode}. On introductory programming exam questions, Codex
scored better than most students \cite{finnieansley2022robots}. Computing
educators have since mapped the opportunities and threats
\cite{becker2023programming,prather2023robots,denny2024genai}; instructors'
plans range from banning the tools to redesigning assessment around them
\cite{lau2023ban}.

Detecting generated text after the fact is unreliable. Weber-Wulff et al.\
tested fourteen detection tools and found none accurate or reliable enough for
academic-integrity decisions \cite{weberwulff2023testing}; Sadasivan et al.\
argued that paraphrasing defeats detectors and that reliable detection becomes
impossible as models approach human text \cite{sadasivan2023detected}; and
Liang et al.\ found that GPT detectors flag essays by non-native English
writers as machine-written far more often than essays by native writers
\cite{liang2023gptdetectors}, a fairness failure with direct consequences for
students. For code, Hoq et al.\ trained classifiers to separate ChatGPT-written
from student-written CS1 submissions \cite{hoq2024chatgpt}, but such models
must be retrained as models and prompts change.

Process data offers a different route: whatever produced the text, a student
who retypes it from another screen types differently from one who composes it,
and a student who pastes it leaves a paste event. This is Nosi's premise
\cite{lueder2026nosi} and the reason this thesis records copying signs
(\cref{sec:weak-labels}). It is also a reason for care: a transcription
measure is a measure of \emph{how} text was typed, not proof of where it came
from.

\section{Fairness of behavioural measures}
\label{sec:rw-fairness}
Biometric systems do not err equally across people. NIST's evaluation of face
recognition found false-match and false-non-match rates that differed across
demographic groups by factors of ten or more for many algorithms
\cite{grother2019frvt}, and commercial face analysis had error rates far
higher for darker-skinned women than for lighter-skinned men
\cite{buolamwini2018gender}. Drozdowski et al.\ surveyed demographic bias
across biometric modalities and found it largely unstudied outside the face
\cite{drozdowski2020bias}. For keystroke dynamics, demographic attributes such
as gender can be predicted from typing features \cite{tsimperidis2018gender},
which means typing carries group information that a detector might pick up,
and typing also differs with skill, keyboard layout and device
\cite{dhakal2018typing}. In education, Baker and Hawn reviewed algorithmic bias
and concluded that most of it is found only when someone looks for it by
subgroup \cite{baker2022bias}. Group fairness criteria, such as equal false
positive rates across groups, are defined in the general literature
\cite{barocas2023fairml}; this thesis reports false alarms and detection by
device class, operating system, language, context and typing speed
(\cref{ch:rq4}) rather than claiming any one criterion is met.

A second, statistical fairness issue is the base rate. When the event being
detected is rare, even a small false-alarm rate means that most alarms are
false: Axelsson made the point for intrusion detection, where it dominates
the usefulness of any detector \cite{axelsson2000baserate}. Takeovers are rare
in any one lab, so precision at a realistic prevalence, not detection rate,
decides what a flag may be used for (\cref{ch:rq4}).

\section{Privacy and ethics of monitoring students}
\label{sec:rw-privacy}
Learning analytics has a developed ethics literature. Slade and Prinsloo
framed it as a moral practice, with students as agents in, not objects of,
the analysis, and data use transparent to them \cite{slade2013ethics}. Pardo
and Siemens set out principles of transparency, student control over data,
security, and accountability \cite{pardo2014ethical}. Drachsler and Greller's
DELICATE checklist turns such principles into questions an institution
answers before deploying analytics, including legitimacy, consent and
anonymisation \cite{drachsler2016delicate}. Prinsloo and Slade also note the
converse duty: an institution that collects data and sees a problem has an
obligation to act on it, which is why what a flag \emph{leads to} must be
designed, not left open \cite{prinsloo2017elephant}. The proctoring studies
above show what happens when these principles are skipped: students experience
monitoring as distrust \cite{balash2021examining,coghlan2021proctor}.

Law sets the outer limits. India's Digital Personal Data Protection Act 2023
requires notice and consent for processing personal data, purpose
limitation and erasure on request \cite{dpdp2023}. The EU AI Act classifies
AI systems used in education to evaluate learning outcomes or to monitor and
detect prohibited behaviour during tests as high-risk, with duties of risk
management, data governance, human oversight and transparency, and it
prohibits emotion recognition in education \cite{euaiact2024}. The EU Act does
not apply in India, but its high-risk requirements are a useful design target
(\cref{sec:ethics}). The editor follows both: it records no camera, screen or
audio, tells students exactly what is recorded, deletes raw key events after
28 days, lets a student erase their data, and gives every decision to a person
(\cref{sec:privacy}).

\section{Summary of the gap}
\label{sec:rw-gap}
\Cref{tab:rw-gap} places the prior work by \emph{when} in a session it checks
authorship and \emph{what} it needs. Identity and proctoring checks act at the
start or watch the person, not the work. Similarity and generated-text
detection act on the finished product. Paste logs and Nosi watch the process,
but for copying, not for a change of author, and Nosi's evidence rests on five
expert participants. Keystroke continuous authentication watches the typist
throughout, but it has been built and evaluated on prose and passwords, with
population-level decisions per window, not on students' code in real classes.
Programming-education work shows that student typing is individual, but does
not monitor a session for a change.

No prior work, to our knowledge, checks \emph{authorship continuity} in the
middle of a programming session from the student's own baseline, estimates
when a change happened, and evaluates this on naturalistic classroom data with
error rates by subgroup and precision at realistic prevalence. That is the
gap this thesis addresses.

\begin{table}[ht]
  \centering
  \caption{Prior work by when it checks authorship and what it needs.}
  \label{tab:rw-gap}
  \small
  \begin{tabularx}{\linewidth}{@{}>{\raggedright\arraybackslash}p{3.3cm}
      >{\raggedright\arraybackslash}p{1.9cm}
      >{\raggedright\arraybackslash}X
      >{\raggedright\arraybackslash}p{2.4cm}@{}}
    \toprule
    Approach & When & Signal & Camera or screen? \\
    \midrule
    Log-in, ID check, proctoring \cite{atoum2017proctoring,nigam2021proctoring}
      & start, whole session & face, room, screen & yes \\
    Code similarity \cite{schleimer2003winnowing,prechelt2002jplag,novak2019similarity}
      & end & finished code across students & no \\
    Generated-text detection \cite{weberwulff2023testing,hoq2024chatgpt}
      & end & finished text or code & no \\
    Paste logs, edit replay
      & throughout & pastes, edit history & no \\
    Nosi \cite{lueder2026nosi}
      & throughout & keystroke rhythm (transcription), population model, $N=5$ & no \\
    Continuous keystroke authentication \cite{bours2012continuous,acien2022typenet}
      & throughout & typing vs.\ enrolled user; prose, passwords & no \\
    Programmer typing \cite{longi2015identification,leinonen2016exams}
      & per sample & identity from typing in courses & no \\
    \midrule
    This thesis
      & throughout & change from the student's own baseline (typing, key pairs, process); estimated switch time; real classes & no \\
    \bottomrule
  \end{tabularx}
\end{table}
